\section{Introducci\'on}
\label{sec:intro}

Un sistema de recuperaci\'on aumentada por generaci\'on (RAG) responde una pregunta con fragmentos recuperados de una colecci\'on y un modelo de lenguaje que redacta a partir de ellos \cite{lewis2020rag}. La promesa es atractiva: actualizar el conocimiento sin reentrenar y citar la evidencia. En la pr\'actica, la calidad depende de una cadena de decisiones que se afectan entre s\'i. Una limpieza agresiva puede borrar la frase que responde; un chunking peque\~no pierde contexto; un codificador entrenado en otro dominio recupera vecinos plausibles pero in\'utiles; un reordenador puede mejorar el orden y a la vez encarecer cada consulta; un prompt permisivo invita a inventar; y sin un control de abstenci\'on el sistema responde con seguridad sobre temas que el corpus no cubre.

La literatura de RAG suele evaluar un componente a la vez. BEIR y MTEB miden recuperadores \cite{thakur2021beir,muennighoff2023mteb}; los benchmarks multiling\"ues como MIRACL comparan codificadores \cite{zhang2023miracl}; RAGAS, ARES, eRAG y RAGChecker eval\'uan la respuesta generada \cite{es2023ragas,saadfalcon2024ares,salemi2024erag,ru2024ragchecker}. Son piezas necesarias, pero no dicen c\'omo se comporta la cadena completa sobre un corpus real, con duplicados, truncamiento, ruido de rastreo y preguntas sin respuesta. Tampoco dicen cu\'anto cuesta cada mejora ni cu\'ando conviene rechazarla.

Este trabajo construye y eval\'ua esa cadena completa sobre un corpus de 169.248 documentos recolectados de prensa colombiana e internacional y Bluesky, 161.868 indexados y 618.286 fragmentos. El sistema corre \'integramente en una estaci\'on de trabajo con una GPU AMD de 16 GB, Milvus como motor vectorial \cite{milvusdocs} y modelos servidos por Ollama \cite{ollamadocs}. No se usan servicios externos de generaci\'on ni de embeddings. La Tabla~\ref{tab:config-final} resume la configuraci\'on final y la Secci\'on~\ref{sec:arquitectura} justifica cada elecci\'on.

\subsection{Preguntas de investigaci\'on}

\begin{enumerate}
\item \textbf{Calidad del corpus.} ?`Qu\'e problemas de limpieza, truncamiento, duplicados e idioma aparecen en un corpus de prensa y redes sociales, y cu\'anto afectan la recuperaci\'on?
\item \textbf{Recuperaci\'on.} ?`Qu\'e combinaci\'on de codificador, \'indice, tama\~no de contexto y reordenamiento maximiza la cobertura y el orden con un coste aceptable?
\item \textbf{Generaci\'on y abstenci\'on.} ?`C\'omo se comportan distintos modelos locales sobre la misma evidencia y qu\'e controles evitan respuestas sin respaldo?
\item \textbf{Evaluaci\'on.} ?`Qu\'e acuerdo alcanzan los anotadores humanos y qu\'e se\~nales de evaluaci\'on autom\'atica son fiables cuando el juez tambi\'en es un modelo local?
\end{enumerate}

\subsection{Aportes}

\begin{enumerate}
\item Un pipeline reproducible de extremo a extremo, con Docker Compose, Milvus, Ollama, FastAPI y Streamlit, que se levanta con un comando y no depende de servicios de pago.
\item Un an\'alisis exploratorio y una textometr\'ia del corpus que cuantifican el texto sin cierre terminal (32,7\,\%, un proxy de truncamiento), ruido de rastreo, duplicados exactos y near-duplicates, y que justifican una limpieza conservadora.
\item Una comparaci\'on pareada de cuatro variantes de limpieza, dos estrategias de segmentaci\'on, tres codificadores, recuperaci\'on densa, dispersa e h\'ibrida, y tres alcances de reordenamiento, con intervalos de confianza por bootstrap.
\item Un gold de 43 preguntas con revisi\'on de tres anotadores y una cuarta pendiente, con acuerdo entre anotadores reportado, y una evaluaci\'on autom\'atica con RAGAS que distingue el ruido del juez de una mejora real.
\item Una compuerta de dominio entrenada con 160 preguntas que rechaza las preguntas sin respuesta con 0 falsos aceptados, y una interfaz que expone la evidencia, las citas y la abstenci\'on.
\end{enumerate}

\subsection{Organizaci\'on}

La Secci\'on~\ref{sec:relacionados} revisa los trabajos relacionados. La Secci\'on~\ref{sec:corpus} describe el corpus y el an\'alisis exploratorio. La Secci\'on~\ref{sec:arquitectura} presenta la arquitectura y justifica cada capa. La Secci\'on~\ref{sec:metodo} detalla el procesamiento y la configuraci\'on. La Secci\'on~\ref{sec:evaluacion} define el protocolo y el gold. La Secci\'on~\ref{sec:experimentos} reporta los experimentos y sus decisiones. Las Secciones~\ref{sec:errores} a \ref{sec:conclusiones} discuten errores, l\'imites y conclusiones.
